\newpage
\section{Diffusion} \label{sec:Diffusion}

Dieses Kapitel beantwortet die Oberfrage und die fünf Unterfragen aus \autoref{sec:einleitung:zielsetzung} und löst die
Ergebnisse von der Fallstudie.

\subsection{Fazit}
\label{sec:diffusion:fazit}

Die Oberfrage verlangte ein Ausführungsmodell, das in einer bestehenden Kubernetes-Plattform
betreibbar bleibt, Mandanten voneinander und aufeinanderfolgende Ausführungen voneinander trennt
und schneller startet als der Ausgangszustand. Die Antwort der Arbeit ist ein Modell mit zwei
Grenzen: Die Mandantengrenze bleibt der Pod unter gVisor, wie im Ausgangszustand, und innerhalb des
Pods erhält jede Ausführung ein eigenes V8-Isolate, das nach dem Lauf verworfen wird. Die Startzeit
nimmt ein Pool vorgehaltener Pods vom Anfragepfad, den eine Control Plane über das vorhandene
Messaging-System zuteilt. Ein Skript erreicht die Außenwelt nur über Bindings, die der Host beim
Erzeugen des Isolates registriert, und eine Kompatibilitätsschicht bildet die Node-Schnittstellen
nach, die der Bestand tatsächlich nutzt.

U1 fragte nach den Anforderungen. Sie stammen aus zwei Quellen der Fallstudie, dem heutigen
Betriebskonzept mit seinen Vertrauensannahmen und der Auswertung von 1.718 Mandanten mit 14.575
Skriptversionen. Der Kern des Problems ist die Verschiebung der Vertrauenseinheit vom Mandanten auf
die einzelne Ausführung, sobald der Code von einem Agenten stammt. Daraus folgen fünf
Muss-Kriterien, die ein Modell einzeln erfüllen muss, und sechs gewichtete Soll-Kriterien
(\autoref{sec:kriterienkatalog}).

U2 fragte nach der Kombination aus Isolationsgrenze, Zuordnungsgranularität und Instanzverwaltung.
Die Muss-Kriterien scheiden Node.js als Vertrauensgrenze, ein prozessweites Berechtigungsmodell, Bun
und workerd aus. Eine nach WebAssembly kompilierte Engine ist zurückgestellt, nicht widerlegt
(\autoref{sec:vorauswahl}). Die Zuordnung ist grob bei der Sicherheitsgrenze, ein Pod je Mandant,
und fein bei den Rechten, feste Bindings je Profil mit einer mandantenweiten Freischaltung für
rohe Sockets. Die Instanzverwaltung hält die Zuteilung vom Anfragepfad fern und entkoppelt die
Poolgröße von der Zahl der Mandanten (\autoref{sec:orchestrierung}).

U3 fragte, ob das Modell Isolation und Zustandsfreiheit einhält. Die neun Prüffälle, die aus dem
Bedrohungsmodell abgeleitet sind, besteht der Prototyp in beiden Profilen alle, der Ausgangszustand sieben nicht
(\autoref{sec:pruefung-isolation}). Zustand überdauert keine Ausführung, ein erschöpfter Heap
trifft nur die verursachende Ausführung, Netz und Dateisystem sind nur über die Bindings
erreichbar, und eine Ausführung kann ihre Identität nicht fälschen. M3 gilt als erfüllt, weil
der Pod in keinem Sicherheitsfeld lockerer ist als der Ausgangszustand (\autoref{tab:m3}).

U4 fragte nach den Messgrößen im Vergleich zum Ausgangszustand. Auf demselben Cluster und unter
gVisor startet der Prototyp mit Pool im Median in 344\,ms statt 2,54 Sekunden und bedient alle 90 statt
10 von 91 Mandanten innerhalb von zwei Sekunden. Ohne Pool liegt er im Median und im Schub auf dem Niveau
des Ausgangszustands und im Tail der verteilten Erstkontakte darüber, der Vorsprung stammt also aus der
Vorhaltung (\autoref{sec:kaltstart}). Der
Overhead im Lastprofil liegt im Median bei 138 statt 259\,ms und im 99.\ Perzentil bei 405\,ms
statt 4,57 Sekunden (\autoref{sec:warm}). Ein Worker-Pod belegt mit 59\,MiB zwei Drittel eines
Runner-Pods, die Steuerung ein Neuntel von Knative Serving, der Pool kostet dauerhaft 1,9\,GiB
(\autoref{sec:durchsatz}). Der Durchsatz liegt bei 500 statt 100 Aufrufen je Sekunde und wird von der Datenbank begrenzt. Beim Knotenentzug
war der Prototyp in acht Läufen höchstens 1,37 Sekunden ohne Ausführungsfähigkeit gegen 1,9 Sekunden
im Ausgangszustand, ohne verlorene Aufträge (\autoref{sec:pruefung-m5}). In der Zusammenführung erreicht der Prototyp auf
vier der sechs Kriterien die Höchstpunktzahl und bei S1 und S3 vier Punkte, bei S1, weil S1b den Zielwert verfehlt, und bei S3, weil der zweite
Prozess den Speicher je Pod erhöht (\autoref{sec:nutzwert}).

U5 beantworten die beiden folgenden Abschnitte. Zuvor gehören die Grenzen der Ergebnisse benannt, die
\autoref{sec:messaufbau} begründet. Hinzu kommen drei Grenzen des Entwurfs. Das Agentenprofil hat bei Kaltstart und
Latenz keinen Vergleichswert, ein Richtlauf zeigt 623\,ms im Median beim ersten Aufruf und 148\,ms warm
(\autoref{sec:kaltstart}). Seine Sicherheitsaussage gilt für aufeinanderfolgende Ausführungen, und nach einem
Ausbruch aus der Engine hält die Mandantengrenze erst, wenn der Auth Callout auf der Plattform eingerichtet
ist (\autoref{sec:betriebsplanung}). Und die Zuverlässigkeit der Zustellung ruht unter höchstens einmaliger
Zustellung auf der Disziplin des Workers (\autoref{sec:zuverlaessigkeit}).

Gegenüber den bekannten Plattformen liegt der Beitrag nicht in einer neuen Isolationstechnik, sondern in der
Kombination für eine Plattform ohne Zugriff auf \ac{KVM}, die über Messaging skaliert, und darin, dass sie gegen
einen produktiven Ausgangszustand gemessen wurde statt gegen generische Vergleichswerte.

\subsection{Designprinzipien und Migrationspfad}
\label{sec:diffusion:prinzipien}

Ein Designprinzip im Sinne von \textcite[S.~1633--1634]{gregor_research_2020} benennt ein Ziel für
einen Nutzer in einem Kontext, die Mechanismen, mit denen ein Implementierer es erreicht, und eine
Begründung dafür, dass die Mechanismen das Ziel erreichen. Für die folgenden Prinzipien ist der
Implementierer der Betreiber einer mandantenfähigen Plattform, der Nutzer der Mandant und die
Ausführungsumgebung der Akteur, der die Mechanismen ausführt. Aufgenommen wird nur, was ein
Schutzziel, einen Mechanismus, eine Kontextbedingung und einen Beleg aus \autoref{sec:Evaluation}
hat. Was nur in dieser Fallstudie gilt oder nicht geprüft wurde, bleibt eine Eigenschaft des
Prototyps.

\begin{description}
  \item[P1 Zwei Grenzen für zwei Schutzziele.] Ziel: Mandanten voneinander und aufeinanderfolgende
    Ausführungen voneinander trennen. Mechanismus: eine vom Kernel durchgesetzte Grenze je Mandant
    und eine Sprachgrenze je Ausführung innerhalb davon. Kontext: fremder Code, dessen
    Vertrauenseinheit die einzelne Ausführung ist. Beleg: Architekturprüfung für die äußere Grenze,
    T1 bis T9 für die innere. Die Sprachgrenze trägt keine Mandantentrennung, die Kernelgrenze
    keine Zustandsfreiheit zu vertretbaren Kosten. Jede Grenze trägt genau das Ziel, für das sie
    billig genug ist.
  \item[P2 Vorhaltung statt Wiederverwendung.] Ziel: Startkosten vom Anfragepfad nehmen, ohne
    Zustand über Ausführungen zu tragen. Mechanismus: ein Vorrat frischer Pods, deren Isolates aus einem
    Snapshot entstehen, dessen Nachfüllen nie auf dem Anfragepfad liegt. Kontext: die Größe des
    Vorrats ist von der Zahl der Mandanten entkoppelt, und es gibt kein Skalieren auf null. Beleg:
    S1a, S1b und der Lauf mit leerem Pool, der zeigt, dass der Vorsprung aus der Vorhaltung
    stammt, nicht aus einem schnelleren Start. Die Wiederverwendung endet an der Grenze, die das
    Schutzziel trägt: Ein Isolate wird nie über Ausführungen, ein Pod nie über Mandanten
    wiederverwendet, und gecacht wird nur Unveränderliches, nie Heap-Zustand.
  \item[P3 Die Fähigkeitsmenge folgt aus der Vertrauensklasse des Codes.] Ziel: einem Skript nur
    die Außenwelt zu öffnen, die seine Herkunft rechtfertigt. Mechanismus: eine feste Menge
    registrierter Bindings je Vertrauensklasse, die die Grenze selbst ist. Die Prüfung der Ziele in
    den Bindings ist die zweite Linie. Kontext: der Inhalt des Codes taugt nicht als Grundlage,
    weil er im Agentenprofil von dem geschrieben wird, dessen Beeinflussbarkeit das
    Bedrohungsmodell beschreibt. Beleg: T7 und T8.
  \item[P4 Der Mandantenkontext liegt außerhalb des Heaps.] Ziel: dass eine Ausführung die Frage,
    in wessen Namen sie handelt, nicht beantworten und damit nicht falsch beantworten kann.
    Mechanismus: Kennung, Profil und zulässige Ziele liegen im Zustand des Hosts, die Bindings
    lesen sie dort, und alles, was das Skript sieht, sind Kopien, die kein Binding liest. Kontext:
    jede eingebettete Engine mit einem Host-Zustand je Instanz. Beleg: T9, in dem jede Kopie
    überschrieben wird und keine Entscheidung sich ändert.
  \item[P5 Ressourcenbegrenzung auf zwei Ebenen.] Ziel: dass eine Ausführung die Ressourcen
    anderer nicht erschöpft und die verursachende Ausführung benannt werden kann. Mechanismus: die
    Mechanismen der Engine, Heap-Limit, zählender Allokator und Watchdog, liefern die Zuordnung und
    einen abfangbaren Fehler. Die Grenze des Betriebssystems je Mandant liefert die Garantie.
    Kontext: Mechanismen der Engine wirken nur an aufgezählten Pfaden. Beleg: T6, in dem der
    Abbruch die verursachende Ausführung trifft und die folgenden fünf unberührt lässt, gegen den
    Prozessabsturz des Ausgangszustands.
\end{description}

Drei Eigenschaften des Prototyps wurden geprüft und nicht aufgenommen. Höchstens einmalige Zustellung folgt aus der
fehlenden Idempotenz von Skriptausführungen und ist eine Eigenschaft der Domäne, keine Entwurfsleistung. Der
Verzicht auf neue Bausteine gehört zum Migrationspfad. Und die Regel, den Weg zu wählen, der die Alternative offen
hält, ist eine Entscheidungsheuristik, die den Ausblick begründet, kein Prinzip für Ausführungsmodelle.

Der Migrationspfad nutzt, dass das Modell keine neuen Bausteine einführt und je Mandant umschaltbar ist: Welcher Stack
die Aufrufe eines Mandanten bedient, entscheidet die Adresse der App in seiner Konfiguration, derselbe Mechanismus,
mit dem \autoref{sec:Evaluation} beide Modelle nacheinander auf demselben Cluster gemessen hat. Die Umstellung erfolgt
schrittweise über ein Feature-Flag, zuerst ein Prozent der Mandanten, dann in Stufen mehr. Nach jeder Stufe werden
Last und Fehlerrate gegen die Werte aus \autoref{sec:Evaluation} geprüft, und eine Stufe wird zurückgenommen, wenn die
Fehlerrate die des Prototyps im Lastprofil übersteigt, durch Umstellen der App-Adresse des Mandanten. Mandanten,
deren Skripte rohe Sockets brauchen, erhalten die mandantenweite Freischaltung vor ihrer Stufe, ein Mandant, dessen
Prüfskripte \texttt{child\_process} ausführen, bleibt bewusst auf dem bisherigen Betrieb. Vor der ersten Stufe sind die
offenen Maßnahmen aus \autoref{sec:betriebsplanung} zu entscheiden, vor allem der Standardwert der
Socket-Freischaltung, der Egress-Default des Agentenprofils und die Ablauffrist auf synchronen Aufträgen, hinzu kommt
die Einrichtung des Auth Callouts durch das \acs{DDP}-Team.

\subsection{Ausblick}
\label{sec:diffusion:ausblick}

Drei Fragen hat die Arbeit geöffnet, ohne sie zu schließen. Die erste betrifft die Engine. Eine nach WebAssembly
kompilierte Engine ist in \autoref{sec:wasm-zurueckgestellt} zurückgestellt und ließe sich später in das gewählte
Modell einsetzen, ohne Grenze und Zuteilung zu ändern. Dringlicher geworden ist die Frage durch die Ankündigung
des Deno-Teams, zu Cloudflare zu wechseln: Die Bindung an \texttt{deno\_core} ist damit eine Abhängigkeit mit
absehbarem Ende, und der Ersatz der Einbettungsschicht ist nach \autoref{sec:betriebsplanung} benennbar, aber nicht
geschätzt.

Die zweite Frage betrifft die Zustellung. Höchstens einmalige Zustellung verliert eine Ausführung beim Absturz
eines Workers still. Ob mindestens einmalige Zustellung mit einer Startmarke je Nachricht
(\autoref{sec:zuteilung}) die Fehlerrate aus \autoref{sec:zuverlaessigkeit} rechtfertigt, ist zu messen, nicht zu
entscheiden.

Die dritte Frage betrifft das Agentenprofil, das nur als Richtwert gemessen ist. Sobald ein Agent Skripte erzeugt,
lassen sich Kaltstart und Overhead für sein Profil erheben. Erst dann zeigt sich, ob die Isolate-Erzeugung bei
kurzen, einmaligen Skripten ins Gewicht fällt und ob sich das vorab erzeugte Isolate aus \autoref{sec:warm}
lohnt.

